\pagestyle{plain}
\chapter*{Resumen}
\addcontentsline{toc}{chapter}{Resumen}

Este Trabajo de Fin de Grado nace de una propuesta del profesor Eduardo
Manuel Sánchez Vila para su grupo de investigación. Aunque al principio fuese algo abstracto, se definieron dos ideas de tfg que juntas construirían un laboratorio completo de estudio de paradigmas de toma de decisiones humanas mediante agentes inteligentes.
Esta memoria recoge y documenta la segunda de estas fases: la infraestructura que
recoge los modelos generados por la primera fase, los ejecuta en un entorno
común y produce en cada simulación artefactos verificables, desde una traza
estructurada del comportamiento hasta un informe comparativo en PDF.

El sistema se organiza como un pipeline de cuatro agentes especializados,
coordinados por un Orchestrator con chat interactivo. El Architect define
el entorno de simulación a partir de una descripción en lenguaje natural;
el Tracker registra eventos, episodios y trayectorias de decisión; el
Analyst identifica patrones de comportamiento y los contrasta con
conocimiento previo; el Reporter sintetiza los hallazgos en un informe
generado mediante LaTeX. El acceso principal es una interfaz web sobre un
backend FastAPI con WebSocket.

La integración con la primera fase del proyecto se realiza por \textit{duck
typing} sobre un contrato mínimo (\texttt{decide}, \texttt{update} y
\texttt{get\_state}), sin acoplamiento de clases entre ambos sistemas. Las
observaciones registradas durante cada simulación se persisten para
alimentar ejecuciones futuras.

\vspace{0.5cm}
\noindent{\bf Palabras clave:} agentes inteligentes, simulación multi-agente,
observación de comportamiento, toma de decisiones, grafos de
conocimiento.
